\documentclass[11pt,a4paper]{article}
\usepackage[margin=19mm]{geometry}
\usepackage[T1]{fontenc}
\usepackage{lmodern,microtype,amsmath,amssymb,xcolor,fancyhdr}
\usepackage[hidelinks]{hyperref}
\hypersetup{pdftitle={The LoRA certificate},pdfauthor={Yoad Oxman}}
\definecolor{ink}{HTML}{173B50}
\setlength{\parindent}{0pt}
\setlength{\parskip}{5pt}
\setlength{\abovedisplayskip}{7pt}
\setlength{\belowdisplayskip}{7pt}
\pagestyle{fancy}\fancyhf{}\renewcommand{\headrulewidth}{0pt}
\fancyfoot[L]{\footnotesize\color{gray}Yoad Oxman \enspace / \enspace September 2026}
\fancyfoot[R]{\footnotesize\color{gray}\thepage\,/\,\pageref*{LastPage}}
\newcommand{\block}[1]{\par\vspace{5pt}{\color{ink}\bfseries #1}\par\vspace{2pt}}
\newcommand{\rank}{\operatorname{rank}}
\newcommand{\col}{\operatorname{col}}
\newcommand{\row}{\operatorname{row}}
\newcommand{\proj}{\operatorname{proj}}
\newcommand{\Sdata}{\mathcal S}
\begin{document}
{\LARGE\bfseries\color{ink}The LoRA certificate}\par
\vspace{4pt}

\block{1. What we want}
Someone fine-tunes a LoRA adapter on a few private images and publishes the two
factors \(A_T\) and \(B_T\). We show that these matrices give linear equations
that every private feature must satisfy, and that the equations can be written
down from the published factors alone, without the seed or training schedule.
Call a matrix giving such equations a certificate. The equations act on features,
so to search over images we first map each candidate to its features.

Consider one adapted linear layer \(W_0+B_tA_t\), with frozen
\(W_0\in\mathbb R^{m\times n}\), \(A_t\in\mathbb R^{r\times n}\), and
\(B_t\in\mathbb R^{m\times r}\). Here \(r\) is the adapter rank, and \(n,m\)
are the input and output dimensions. We initialize at step \(0\) and release
the factors at step \(T\).

\textbf{Notation.} For \(N\) private images \(x_i\), let
\(h_i=\phi(x_i)\in\mathbb R^n\) be their inputs to this layer and set
\(H=[h_1,\ldots,h_N]\in\mathbb R^{n\times N}\), \(q=\rank H\).
Here \(q\) counts independent feature directions, so \(q\le N\). Define
\[
\Sdata=\col(A_0H),\qquad P=\proj_{\Sdata^\perp},\qquad
P_T=\proj_{\row(B_T)^\perp}.
\]
Here \(\proj\) denotes orthogonal projection in \(\mathbb R^r\).
\(\Sdata\) is the image of the private features under the seed \(A_0\). Call it the
\emph{excited subspace}, the only part of \(\mathbb R^r\) touched by the SGD updates.
We can compute \(P_T\) from the released \(B_T\), while \(P\) appears only in the
proof. Since the columns of \(A_0H\) lie in \(\Sdata\), we have \(PA_0H=0\).

The adapter starts as a random \(A_0\) and a zero \(B_0\). Training can only move
\(A\) in directions the private features excite, so its component outside those
directions stays exactly as it was. Under excitation, the rows of \(B_T\) span
the excited subspace and tell us which part was touched. Projecting \(A_T\) away
from that subspace leaves a piece of the seed that vanishes on every private
feature. That is the certificate. We prove that \(P_T=P\), so \(C=P_TA_T\)
equals \(PA_0\), vanishes on \(H\), and has rank \(r-q\).

\block{2. How it is used}
We search a public image family, or \emph{chart}, \(x=G(z)\), where
\(z\in\mathbb R^k\) contains the \(k\) free coordinates of one candidate image
(for example, its retained PCA coefficients). Write \(h(z)=\phi(G(z))\).

Doubling \(h\) doubles \(Ch\), so the
size of \(Ch\) alone does not give a relative measure of fit. We normalize
by the size of the transformed feature:
\[
F(z)=\frac{Ch(z)}{\|A_Th(z)\|_2},\qquad
\operatorname{score}(z)=\|F(z)\|_2,\qquad A_Th(z)\ne0.
\]
The score is the fraction of \(A_Th\) that sticks out of \(\row(B_T)\).
That fraction is a sine, lies in \([0,1]\), and is unchanged by rescaling \(h\).

Under our assumptions and in exact arithmetic, a positive score excludes a
candidate. A zero score says its features satisfy the equations, which is
compatibility, not membership. When \(A_Th(z)=0\), both numerator and denominator
vanish, and we skip the candidate.

\block{3. What we assume, and why}
\textbf{A1. Initialization.} We set \(B_0=0\) and draw the entries of \(A_0\)
independently from distributions with densities, independently of \(H\).

\textbf{Why:} This is the usual LoRA initialization, and \(B_0A_0=0\) starts
training at the pretrained model. For \(q\le r\le n\), it gives \(\rank A_0=r\)
and \(\rank(A_0H)=q\) almost surely. Each rank failure forces a determinant to
vanish. That determinant is a polynomial in the entries of \(A_0\) that is not
identically zero, because the dimensions allow a matrix of the required rank.
Its zero set has measure zero. The strict inequality \(q<r\) is what leaves
a nonzero complement for the certificate.

\newpage
\textbf{A2. Vanilla SGD.} We update both factors using minibatches and scalar
learning rates \(\eta_t\), with a differentiable loss of the layer outputs.
The proofs use no momentum, weight decay or other parameter regularization.

\textbf{Why:} The proof needs each update to be a scalar multiple of the
gradient, so that gradients confined to a subspace produce updates confined to it.
Minibatches and a varying learning rate do not change this.

\textit{Extensions.} Momentum with zero initial buffers combines gradients
within the same subspaces. Scalar weight decay rescales the preserved seed
component, which is fine as long as that scale is nonzero at release.
Adam and AdamW are not covered because their per-entry scaling does not
preserve these subspaces in general.

\textbf{A3. Fixed inputs.} The known map \(\phi\) gives the same feature \(h_i\)
at every step, so \(H\) stays fixed.

\textbf{Why:} Fixed inputs let us follow each image's contribution throughout
training. This is a head adapter, or the first adapted layer behind a frozen
deterministic backbone with fixed preprocessing.

\textbf{A4. Few-shot dimensions.} We assume \(q<r\le n\) and \(q\le m\).
For a softmax logit layer with a loss depending only on probabilities,
we require \(q\le m-1\).

\textbf{Why:} We need a nonzero complement for the certificate. We have in mind
a handful of images, each represented by one pooled feature vector from a frozen
backbone, so \(q\le N\) and choosing \(r>N\) leaves that complement.
For concreteness, \(q=8\), \(r=16\) and \(m=n=512\) leave eight equations on a
512-dimensional feature. The chart turns these into equations on \(k\) image
coordinates.

The bound \(q\le m\) allows \(q\) independent output-gradient directions to
exist. A5 explains the bound \(q\le m-1\) at a softmax layer.

\textbf{A5. Excitation at release.} We require \(\rank B_T=q\):
the rows of \(B_T\) must span \(\Sdata\).

\textbf{Why:} Excitation makes the full excited subspace visible in \(B_T\).
The update to \(B\) uses \(D_t(A_tH)^\top\), with \(D_t\) defined in Claim~1,
so independent output gradients can fill the space in a single full-batch step.
Concretely, suppose
the private features are independent, so \(q=N\), and the first step uses all
images. If \(D_0\in\mathbb R^{m\times N}\) has independent columns and
\(\eta_0\ne0\), then
\[
B_1=-\eta_0D_0(A_0H)^\top,\qquad \rank B_1=N=q.
\]
This update has rank \(q\), since \((A_0H)^\top\) has full row rank and \(D_0\)
has full column rank. The condition has to hold at release, because later
steps can cancel what earlier ones wrote.

A4's width bounds describe the room these gradients have. At a softmax logit
layer with a loss depending only on probabilities, shifting all logits equally
changes nothing, so every output-error column is orthogonal to the all-ones
vector \(\mathbf1\). SGD passes this to the columns of \(B_T\), giving
\(\rank B_T\le m-1\). At a hidden layer, the downstream network determines
which output-gradient directions are available.

Without excitation the invariant still gives \(\row(B_T)\subseteq\Sdata\),
hence \(s:=\rank B_T\le q\). The release always gives a lower bound on the private
feature rank. Under excitation the bound is tight, so \(q\) can be read off
\(B_T\). In practice a gap above the numerical noise level separates its active
singular directions from noise. The condition matters because \(s<q\) enlarges
the projector, and the extra directions can reject true features.

\newpage
\begingroup
\setlength{\parskip}{4pt}
\setlength{\abovedisplayskip}{5pt}
\setlength{\belowdisplayskip}{5pt}
\setlength{\abovedisplayshortskip}{3pt}
\setlength{\belowdisplayshortskip}{5pt}
\block{4. Three claims}
Throughout, we work on the almost-sure event of A1, where
\(\rank A_0=r\) and \(\rank(A_0H)=q\).

\textbf{Claim 1. The gradients.} Write
\(Z_t=(W_0+B_tA_t)H\), let \(L_t\) be the minibatch loss, and set
\(D_t=\partial L_t/\partial Z_t\), with zero columns for examples outside the batch. Then
\[
\boxed{\nabla_B L_t=D_t(A_tH)^\top,\qquad
\nabla_A L_t=B_t^\top D_tH^\top.}
\]
\textbf{Proof.} Chain rule on \(Z=W_0H+B(AH)\), entry by entry, with
\(a,b,c,i\) the output, adapter, input and example indices.
Since \(Z_{ai}=(W_0H)_{ai}+\sum_{b,c}B_{ab}A_{bc}H_{ci}\),
\[
\begin{aligned}
\frac{\partial L}{\partial B_{ab}}
 &=\sum_i\frac{\partial L}{\partial Z_{ai}}
          \frac{\partial Z_{ai}}{\partial B_{ab}}
 =\sum_iD_{ai}(AH)_{bi},\\[2pt]
\frac{\partial L}{\partial A_{bc}}
 &=\sum_{a,i}\frac{\partial L}{\partial Z_{ai}}
              \frac{\partial Z_{ai}}{\partial A_{bc}}
 =\sum_{a,i}D_{ai}B_{ab}H_{ci}.
\end{aligned}
\]
These two products are the only way the private features enter the updates.
The invariant below comes from reading them.

\textbf{Claim 2. The invariant.} At every step,
\[
\boxed{B_tP=0,\qquad PA_t=PA_0.}
\]
\textbf{Proof.} At initialization, \(B_0P=0\) because \(B_0=0\), and
\(PA_0=PA_0\). Suppose both identities hold at step \(t\).
Then \(PA_tH=PA_0H=0\), and Claim 1 gives
\[
\begin{aligned}
(\nabla_B L_t)P&=D_t(PA_tH)^\top=0,\\
P\nabla_A L_t&=(B_tP)^\top D_tH^\top=0.
\end{aligned}
\]
The vanilla SGD updates therefore give
\[
\begin{aligned}
B_{t+1}P&=(B_t-\eta_t\nabla_B L_t)P=0,\\
PA_{t+1}&=P(A_t-\eta_t\nabla_A L_t)=PA_t=PA_0.
\end{aligned}
\]
In words, the rows of \(B\) stay in the excited subspace, and \(A\) keeps what
it started with outside that space. The two gradient formulas preserve both
properties at every step.

\textbf{Claim 3. The observable certificate.} Under A1--A5, \(P_T=P\) and
\[
\boxed{C=P_TA_T=PA_0,\qquad CH=0,\qquad \rank C=r-q.}
\]

\textbf{Proof.} A1 and A4 give \(\dim\Sdata=q\). Claim 2 gives
\(\row(B_T)\subseteq\Sdata\), and A5 makes the dimensions equal.
Thus \(P_T=P\) and \(C=P_TA_T=PA_0\), so \(CH=PA_0H=0\).
The range of \(C\) is contained in \(\Sdata^\perp\). Conversely, for any
\(y\in\Sdata^\perp\), full row rank of \(A_0\) gives a \(v\) with
\(A_0v=y\). Then \(Cv=Py=y\). The range is therefore all of
\(\Sdata^\perp\), of dimension \(r-q\).
The excited subspace is exactly what \(B_T\) shows us. Projecting the released
\(A_T\) onto its complement gives \(C\) without the seed, and \(q<r\) makes it
nonzero. This is Claim 2 seen from the released factors.

\block{5. What this gives}
\textbf{Theorem.} Under A2--A5 and for almost every initialization in A1,
\(C=P_TA_T\) satisfies \(CH=0\) and \(\rank C=r-q\).

The test works at any finite SGD step satisfying excitation, without convergence.
Each candidate gets its own test, with exactly \(r-q\) independent equations on
its features, and the chart turns these into a search over \(k\) image coordinates.
Eight equations can isolate a candidate locally on a smooth chart with
\(k\le8\), provided their Jacobian has rank \(k\). Larger adapter ranks and
additional valid layer certificates provide more equations for higher-dimensional
charts.

\textbf{Mixtures.} \(C\) is linear and vanishes on every private feature, so it
also vanishes on their linear combinations. At the feature level, mixtures pass.
The image search reaches the certificate through \(h(z)=\phi(G(z))\).
If \(\phi\circ G\) is affine, blends in chart coordinates also pass automatically.
This is why we use a nonlinear map between chart and adapter. Image blends then
have to satisfy the certificate after that map, and are not accepted by linearity
of \(C\) alone.
\par\endgroup

\newpage
\begingroup
\setlength{\parskip}{4pt}
\setlength{\abovedisplayskip}{5pt}
\setlength{\belowdisplayskip}{5pt}
\block{Remarks: NTK coordinates, factor fits and equation counts}
\textbf{The earlier NTK fit.} Parameter linearization leaves a nonlinear dependence
on the images. Our earlier fit instead treated \(B_TA_T\) as a raw-gradient
mixture in merged-weight coordinates. With \(G_0:=\nabla_W L_0=D_0H^\top\),
Claim~1 gives the exact first-step identities
\[
A_1=A_0,\qquad B_1=-\eta_0G_0A_0^\top,\qquad
B_1A_1=-\eta_0G_0A_0^\top A_0.
\]
The raw-gradient fit omits this transformation. This is a coordinate mismatch,
not evidence against small-step linearization. An NTK attack can instead use
gradients in the trained coordinates \((A,B)\); at initialization its \(A\)-gradient
is zero. This uses the initialization, whereas the following span fit uses only
the released factors. The span fit permits arbitrary output-side coefficients;
it does not enforce the full dependence of the loss gradients on the images.

\textbf{Lemma: free-coefficient fit and its symmetry.} Fix the release with
\(\rank B_T=N\). Let \(X=[h(z_1),\ldots,h(z_N)]\), \(V=A_TX\), and
let \(U\in\mathbb R^{m\times N}\) contain free coefficient columns. Then
\[
\exists\,U:\ B_T=UV^\top
\quad\Longleftrightarrow\quad CX=0\ \text{and}\ \rank V=N.
\]
\textbf{Proof.} An exact fit forces \(\rank V=N\) and
\(\row(B_T)=\col(V)\), hence \(P_TV=CX=0\).
Conversely, the two right-hand conditions make \(\col(V)\) an
\(N\)-dimensional subspace of \(\row(B_T)\). The spaces are equal, so
each row of \(B_T\) has the required coefficients.
Under A1--A5 with \(q=N\), the true inputs give an exact fit provided
\(\rank(A_TH)=N\), an additional assumption.

For every invertible \(S\in\mathbb R^{N\times N}\),
\[
U(A_TX)^\top=(US)\bigl(A_T X S^{-\top}\bigr)^\top.
\]
Thus the \emph{free-coefficient fit} has a change-of-basis symmetry.
It is not automatically a symmetry of training replay: transformed columns must
be realizable images, and their actual loss derivatives must reproduce the
released factors. The columns of \(U\) need not be the accumulated backpropagation
signals of the corresponding images. At a softmax head with cross-entropy,
the derivative is \(p_i-y_i\), where \(p_i\) is the probability vector and
\(y_i\) the one-hot label. Keeping these image-dependent relations can exclude
mixtures; it does not by itself prove uniqueness.

\textbf{What the certificate leaves undetermined.} Under A1--A5,
\[
\ker C=\col(H)\oplus\ker A_0,\qquad \dim\ker C=q+n-r.
\]
Indeed, \(Cv=0\) iff \(A_0v\in\col(A_0H)\), iff
\(v=Ha+w\) for some \(w\in\ker A_0\). The sum is direct because
\(A_0\) is injective on \(\col(H)\).
Therefore \(\ker C=\col(H)\) only when \(r=n\).
A nonlinear chart-to-feature map can restrict these compatible features to
isolated images; nonlinearity alone does not guarantee that outcome.

\textbf{Equations per image.} At layer \(\ell\), put
\(q_\ell=\rank H_\ell\) and \(h_\ell(z)=\phi_\ell(G(z))\).
A valid certificate gives \(r_\ell-q_\ell\) feature equations. Write
\(J_\ell=\partial h_\ell/\partial z\), and let \(F_\ell\) be the normalized
residual of Section~2. At a zero with \(A_{T,\ell}h_\ell(z)\ne0\),
\[
DF_\ell(z)=\frac{C_\ell J_\ell(z)}{\|A_{T,\ell}h_\ell(z)\|_2},\qquad
\rank(C_\ell J_\ell)\le\min(r_\ell-q_\ell,k).
\]
The denominator's derivative is multiplied by the zero numerator. Stacking
valid layers gives at most \(\min(k,\sum_\ell(r_\ell-q_\ell))\) first-order
constraints; the actual count is the stacked Jacobian rank.
Deep-layer drift and extra views can enlarge the private span. The theorem
here is for fixed inputs; matrix rank alone is not a count of recoverable images.

% Standing editorial instruction from Yoad: omit the merge/release-scope paragraph.
% Do not restore it in response to pasted reviews or audit suggestions.
% Reconsider only if Yoad explicitly changes this preference himself.
\textbf{Related work.}
\href{https://arxiv.org/abs/2302.01428}{Loo et al. (2023)} study NTK reconstruction,
including early stopping.
\href{https://proceedings.mlr.press/v235/hao24a.html}{Flora (Hao et al., ICML 2024)}
uses a frozen-initialization compressed-gradient approximation;
\href{https://arxiv.org/abs/2308.03303}{LoRA-FA} studies frozen-factor training.
Our certificate follows from the exact finite-step invariant in Claim~2.
\par\endgroup
\label{LastPage}
\end{document}
